\documentclass[10pt]{article}
\usepackage[margin=0.62in]{geometry}
\usepackage{enumitem}
\usepackage{hyperref}
\usepackage{titlesec}
\pagenumbering{gobble}
\setlist[itemize]{leftmargin=*,nosep}
\titleformat{\section}{\large\bfseries}{}{0pt}{}[\titlerule]
\begin{document}
\begin{center}
{\LARGE Srikrishnaa Vadivel}\\
Machine Learning Engineer \textbar{} Scientific ML \textbar{} Equivariant Models\\
\href{mailto:srikrishnaacareers@gmail.com}{srikrishnaacareers@gmail.com} \textbar{} 516-853-5663\\
\href{https://github.com/krishnaa423}{github.com/krishnaa423} \textbar{}
\href{https://leetcode.com/u/krishnaa42342/}{leetcode.com/u/krishnaa42342} \textbar{}
\href{https://hub.docker.com/u/krishnaa42342}{hub.docker.com/u/krishnaa42342}
\end{center}
\section{Summary}
PhD researcher applying machine learning to first-principles materials and
molecular systems. Experience spans PyTorch, equivariant graph networks,
diffusion models, distributed training concepts, HPC data pipelines, and
physics-based loss design.
\section{ML Experience}
\textbf{Yale University, PhD Researcher} \hfill 2021--present
\begin{itemize}
\item Developed first-principles datasets and workflows for self-trapped
exciton calculations on leadership-class HPC resources.
\item Prototyped EGNN-style force models with energy gradients
$F_{ai}=-\partial E/\partial R_{ai}$ and molecular graph construction.
\item Built diffusion-model notes and code scaffolds covering DDPM training,
conditional inference, transformers, and index-notation derivations.
\end{itemize}
\textbf{Probe Technologies / Weatherford} \hfill 2019--2021
\begin{itemize}
\item Built GPU-accelerated processing and visualization for large acoustic
waveform data streams using CUDA and OpenGL.
\item Engineered production data tooling around ASP services and Microsoft
database systems.
\end{itemize}
\section{Selected Projects}
\begin{itemize}
\item \textbf{EGNN QM9 Force Model}: PyTorch message-passing model, force loss,
synthetic graph fallback, and training-curve generation.
\item \textbf{Diffusion Notes}: DDPM equations with $\alpha_t$, $\beta_t$,
conditional training, and inference in index notation.
\item \textbf{LiF Band Scaffold}: reciprocal-space Hamiltonian assembly and
eigenvalue plotting for physics-informed ML context.
\end{itemize}
\section{Education}
Yale University, PhD Electrical Engineering, 2021--present\\
Yale University, MPhil Electrical Engineering, 2023\\
Yale University, MS Electrical Engineering, 2023\\
Cornell University, MEng ECE, 2017--2018\\
Cornell University, BA ECE, 2013--2017
\section{Skills}
PyTorch, NumPy, SciPy, pandas, scikit-learn, Hugging Face Transformers,
Accelerate, datasets, CUDA, MPI, OpenMP, Linux, Python, C++, Docker,
Kubernetes, PETSc, SLEPc.
\end{document}
